% Zone „Das kennst du schon“ – aus bank/rotationsvolumen/zone.jsonl (zusammenbau v0.1)
\einheitenkopf[zone][Kennst du schon]{Das kennst du schon}
\setcounter{aufgabe}{0}

%% TODO Titel: Ich-kann-Satz fehlt in der Bank (Platzhalter: Kettenname) – Nr. 1
%% TODO Anweisung: ein Satz über der ersten Teilaufgabe fehlt in der Bank – Nr. 1
\begin{aufgabe}{Binomisch quadrieren}
\begin{teile}
\teil $(x + 3)^2$ ausmultipliziert? \leerfeld
\teil $(0{,}5x + 2)^2$ ausmultipliziert? \leerfeld
\end{teile}
\end{aufgabe}

%% TODO Titel: Ich-kann-Satz fehlt in der Bank (Platzhalter: Kettenname) – Nr. 2
%% TODO Anweisung: ein Satz über der ersten Teilaufgabe fehlt in der Bank – Nr. 2
\begin{aufgabe}{Bestimmte Integrale berechnen}
\begin{teile}
\teil Integral von 0 bis 2 über $3x^2\,dx$? \leerfeld
\teil Integral von −2 bis 1 über $x^2\,dx$? \leerfeld
\end{teile}
\end{aufgabe}

%% TODO Titel: Ich-kann-Satz fehlt in der Bank (Platzhalter: Kettenname) – Nr. 3
%% TODO Anweisung: ein Satz über der ersten Teilaufgabe fehlt in der Bank – Nr. 3
\begin{aufgabe}{Zylinder-, Kugel- und Prismenvolumen aus den Sek-I-Formeln}
\begin{teile}
\teil Zylinder mit $r = 3$ cm und $h = 5$ cm – Volumen? V = \leerfeld[cm³]
\teil Kugel mit dem Durchmesser 5 cm – Volumen? V = \leerfeld[cm³]
\end{teile}
\end{aufgabe}

%% TODO Titel: Ich-kann-Satz fehlt in der Bank (Platzhalter: Kettenname) – Nr. 4
%% TODO Anweisung: ein Satz über der ersten Teilaufgabe fehlt in der Bank – Nr. 4
\begin{aufgabe}{Volumeneinheiten und Dichte}
\begin{teile}
\teil 2,5 dm³ – wie viele Liter, wie viele cm³? \leerfeld[l] = \leerfeld[cm³]
\teil 0,4 m³ Erde, Dichte 0,8 t/m³ – Masse in kg? m = \leerfeld[kg]
\end{teile}
\end{aufgabe}

%% TODO Titel: Ich-kann-Satz fehlt in der Bank (Platzhalter: Kettenname) – Nr. 5
%% TODO Anweisung: ein Satz über der ersten Teilaufgabe fehlt in der Bank – Nr. 5
\begin{aufgabe}{Satz des Pythagoras am Kreisschnitt}
\begin{teile}
\teil Kugel mit Radius 5 cm, ebener Schnitt 3 cm vom Mittelpunkt – Radius des Schnittkreises? r = \leerfeld[cm]
\teil Kugel mit Radius 2,5 dm, ebener Schnitt 1,5 dm vom Mittelpunkt – Radius des Schnittkreises? r = \leerfeld[dm]
\end{teile}
\end{aufgabe}

%% TODO Titel: Ich-kann-Satz fehlt in der Bank (Platzhalter: Kettenname) – Nr. 6
%% TODO Anweisung: ein Satz über der ersten Teilaufgabe fehlt in der Bank – Nr. 6
\begin{aufgabe}{Binomisch quadrieren – Fehler finden}
\begin{teile}
\teil Finde den Fehler. \rechnung{(x + 5)^2 &= x^2 + 25}
\end{teile}
\end{aufgabe}

%% TODO Titel: Ich-kann-Satz fehlt in der Bank (Platzhalter: Kettenname) – Nr. 7
%% TODO Anweisung: ein Satz über der ersten Teilaufgabe fehlt in der Bank – Nr. 7
\begin{aufgabe}{Binomisch quadrieren – selbst rechnen}
\begin{teile}
\teil $(x + 6)^2$ ausmultipliziert? \leerfeld
\end{teile}
\end{aufgabe}
